\section{Interpretación y amenazas a la validez}
\subsection{Alcance}
Q1/A/C1/R1 no permite generalizar a otros perfiles, concurrencias o
consultas. Las operaciones comparten procesos y cachés, de modo que no
reemplazan repeticiones independientes.

Promediar percentiles no produce un percentil global. La inversa de QPS
no sustituye necesariamente la latencia media, por los límites y
criterios de contabilización.

\subsection{CPU}
PostgreSQL comunica 0.083895 segundos de throttling y MongoDB cero.
La señal positiva no demuestra que CPU explique la diferencia temporal.
Cero no excluye espera de E/S u otros cuellos de botella.

La cuota define capacidad permitida y no uso constante. Sin B/C ni
C10/C50 no se caracteriza respuesta a cuotas menores o saturación.

\subsection{Memoria}
RSS exige revisar alcance y unidades. No representa automáticamente toda
la memoria atribuible: memoria compartida, procesos y caché del host
pueden tener coberturas diferentes.

PostgreSQL distingue buffers compartidos de memoria de operaciones.
\texttt{work\_mem} no es un límite global y puede multiplicarse con
sesiones y operaciones \cite{pgmemory}. WiredTiger mantiene caché propia
e interactúa con la caché del sistema de archivos \cite{wiredtiger}.
Estas explicaciones no atribuyen por sí solas el RSS de las ventanas.

No se afirman tamaños internos sin configuración ni se suman señales
heterogéneas para fabricar una memoria total.

\subsection{Spill}
Las políticas y planes diferentes son explicaciones posibles del spill,
pero no demuestran causalidad. Se requiere evidencia del operador
y de la ventana para cuantificar trabajo temporal.

El permiso de disco de MongoDB condiciona etapas, sin obligarlas a
derramar \cite{mongolimits}. No equivale a \texttt{work\_mem}.
Las variantes se analizan separadamente.

\subsection{Índices y semántica}
La aceptación del responsable no transforma un índice textual en temporal.
La elegibilidad Q1 no demuestra la ruta de Q2. La equivalencia lógica
no obliga a estructuras internas iguales.

\subsection{Amenazas}
\begin{itemize}
\item \textbf{Interna:} orden, cachés, carga, límites, ventanas y resets.
\item \textbf{Constructo:} unidad RSS, spill, cálculo de QPS y errores.
\item \textbf{Comparación:} variantes, índices, versiones y planes.
\item \textbf{Externa:} un host, una semilla y dos consultas.
\item \textbf{Conclusión:} una repetición por motor no estima variabilidad.
\end{itemize}
Las 135 pruebas del preflight no son medición de ahorro temporal.
El bloqueo al cambiar de motor permanece como incidencia operativa.
